\bookchapter{บทที่ 12: ภาพยนตร์เพลงบันเทิงกับการอ่าน}

\chapterintro{มือถือตอนนี้แทนที่ MP3 รีโมตทีวี Kindle และเครื่องเกมพกพาไปเกือบหมดแล้ว แต่พูดตรงๆ หลายคนใช้แค่ฟังก์ชันพื้นฐานของแอปพวกนี้ ยังมีของดีๆ ที่ใช้ได้อีกเยอะที่ยังไม่เจอ บทนี้มาคุยกันว่าทำยังไงให้ความบันเทิงบนมือถือลื่นและสบายขึ้น}

\booksection{หนึ่ง ฟังเพลง: เลือกแอปไหน ฟังยังไงให้สบาย}

แอปเพลงตอนนี้ฟังก์ชันใกล้เคียงกันหมด แต่แต่ละตัวมีจุดเด่นของตัวเอง เลือกถูกช่วยให้สบายใจขึ้นเยอะ:

\begin{itemize}
  \item \textbf{NetEase Cloud Music}: จุดแข็งคือชุมชนเพลย์ลิสต์กับคอมเมนต์ คุณอยากฟัง "เพลงเหมาะกับวันฝนตก" หรือ "เพลงบรรเลงไว้ผ่อนตอนเรียนเหนื่อย" ค้นหาแป๊บเดียวมีคนจัดเพลย์ลิสต์ให้แล้ว คอมเมนต์บางทีก็เด็ดกว่าเพลงอีก ข้อเสียคือลิขสิทธิ์บางเพลงไม่ครบ ฟังๆ ไปเพลงเทาไปบางเพลง
  \item \textbf{QQ Music}: ลิขสิทธิ์ครบที่สุด เพลงดังอย่าง Jay Chou หรือ Mayday แทบมีแค่ QQ Music เดียว ตัวเลือกคุณภาพเสียงก็อุดม สมาชิกเปิดเสียงไร้สูญเสียได้ ถ้าคุณสำคัญเรื่อง "อยากฟังเพลงไหนก็ค้นเจอ" QQ Music คือทางเลือกที่มั่นคงที่สุด
  \item \textbf{Apple Music}: ตัวเลือกแรกของผู้ใช้แอปเปิล อินเทอร์เฟซสะอาด ไม่มีโฆษณา คุณภาพเสียงดี รองรับเสียงไร้สูญเสียและ Dolby Atmos แต่อัลกอริทึมแนะนำไม่เข้าใจรสนิยมผู้ใช้ภาษาจีนเท่าแอปในประเทศ วัฒนธรรมเพลย์ลิสต์ก็ไม่เจริญเท่า NetEase เหมาะกับคนที่ชอบประสบการณ์เรียบง่าย ไม่สนเรื่องโซเชียลมากนัก
  \item \textbf{Spotify}: ผู้ใช้ทั่วโลกเยอะที่สุด อัลกอริทึมแนะนำได้รับการยอมรับว่าแข็งแรงที่สุด "Discover Weekly" กับ "รายวัน" แม่นยำน่ากลัว ใช้นานเท่าไรยิ่งเข้าใจคุณ ข้อเสียคือเวอร์ชันฟรีโฆษณาเยอะ และลิขสิทธิ์เพลงจีนบางเพลงไม่ครบ ต้องใช้เครือข่ายพิเศษถึงจะเข้าได้
\end{itemize}

\textbf{เคล็ดลับฟังเพลงสักหน่อย}: ถ้าใช้หูฟังบลูทูธ เข้าไปดูในการตั้งค่าแอปเพลงว่ามี “ปรับเอฟเฟกต์เสียง” หรือ “อีควอไลเซอร์” ไหม ปรับเบสกับเสียงแหลมนิดหน่อย เอฟเฟกต์ดีขึ้นชัด อีกอย่างการฟังเพลงสตรีมมิงกินดาต้ามาก อย่าลืมเปิด “เล่นเสียงไร้สูญเสียเฉพาะใต้ Wi-Fi” ในการตั้งค่า หรือดาวน์โหลดเพลย์ลิสต์โปรดเตรียมไว้ฟังออฟไลน์

\booksection{สอง พอดแคสต์: “เพื่อนคู่หู” ระหว่างเดินทางไปทำงาน}

พอดแคสต์ คนที่ไม่เคยฟังจะคิดว่าไม่เกี่ยวกับตัวเอง พอเริ่มฟังปุ๊บก็ไปไม่เป็นแล้ว แท้จริงคือ “เน็ตฟลิกซ์เวอร์ชันรายการวิทยุ” — มีคนคุยเรื่องเทค มีคนเล่าประวัติศาสตร์ มีคนอ่านเรื่องเล่า อยากฟังแนวไหนก็มี

\begin{itemize}
  \item \textbf{แอป Xiaoyuzhou}: แอปพอดแคสต์ที่ทำดีที่สุดในจีน อินเทอร์เฟซสะอาด ไม่มีฟังก์ชันอื้อฉ่าว "แนะนำประจำวัน" และ "สรรหาโดยบรรณาธิการ" ช่วยให้เจอรายการใหม่ คอมเมนต์ก็คึกคัก ที่ฉันฟังเยอะที่สุดคือ "Huzuohuyou" (วัฒนธรรมประวัติศาสตร์) กับ "Buheshiyi" (ประเด็นสังคม) แต่ละตอนเหมือนนั่งฟังกลุ่มคนเก่ง ๆ คุยกัน
  \item \textbf{Apple Podcasts}: มาในตัว iPhone ไม่ต้องดาวน์โหลดแอปเพิ่ม เชื่อมกับ CarPlay ได้ดีที่สุด สะดวกมากตอนขับรถ
  \item \textbf{เทคนิคการฟัง}: แอปพอดแคสต์ส่วนใหญ่รองรับความเร็ว 1.2 เท่าหรือ 1.5 เท่า คุ้นแล้วประหยัดเวลาได้เยอะ ฉันปกติฟังภาษาจีนที่ 1.2 เท่า ภาษาอังกฤษที่ 1.5 เท่า อีกอย่างต้องเปิด "ดาวน์โหลดตอนใหม่อัตโนมัติ" กับ "ลบเมื่อฟังจบอัตโนมัติ" มือถือจะได้ไม่ถูกไฟล์เสียงที่ฟังแล้วกองเต็ม
\end{itemize}

\booksection{สาม ดูวิดีโอ: วิธีเปิดที่ถูกต้องจากเครื่องเล่นไปสตรีมมิง}

หลายคนไม่รู้ว่าฟอร์แมตที่มือถือเล่นได้เยอะกว่าที่แอปสตรีมมิงมีให้มากๆ เครื่องเล่นในเครื่องดีๆ แก้ปัญหาได้เยอะ

\begin{itemize}
  \item \textbf{เครื่องเล่น VLC}: โอเพนซอร์สฟรี เกือบจะเล่นได้ทุกฟอร์แมต — MKV, AVI, FLAC, แม้แต่อิมเมจแผ่น ISO ถ้าคุณดาวน์โหลดวิดีโอที่เปิดได้บนคอมแต่เปิดไม่ได้บนมือถือ แปดทบสิบคือปัญหาฟอร์แมต ติดตั้ง VLC ก็จบ มันรองรับซับไตเติลภายนอกด้วย เอาไฟล์ซับ .rt กับวิดีโอใส่โฟลเดอร์เดียวกัน VLC จะโหลดให้เอง
  \item \textbf{nPlayer}: ถ้าอยากได้ประสบการณ์ประณีตขึ้น แอปเสียเงินนี้คุ้มที่จะซื้อ อินเทอร์เฟซสวย ท่าทางใช้ลื่นมือ — ปัดขึ้นลงฝั่งซ้ายปรับความสว่าง ฝั่งขวาปรับเสียง ปัดซ้ายขวาเดินหน้าถอยหลัง รองรับ DTS และเสียง Dolby ดูหนังใหญ่ได้อารมณ์สมจริงมากขึ้น
  \item \textbf{Bilibili (สถานี B)}: ไม่ใช่แค่เว็บวัยรุ่นสายสองมิติแล้ว ตอนนี้ทรัพยากรการเรียนรู้บนนี้เยอะจนเว่อร์ สอบเข้ามหาวิทยาลัย เขียนโค้ด ถ่ายภาพ ทำอาหาร สอนอะไรก็มี วัฒนธรรมดันมูคือวิญญาณของ Bilibili ดูวิดีโอไปเปิดดันมูไป เหมือนมีคนกลุ่มหนึ่งนั่งดูเคียงคุณ แต่ดูวิดีโอการเรียนแนะนำปิดดันมู ไม่งั้นเสียสมาธิง่าย
  \item \textbf{แคสต์ไปทีวี}: จอมือถือใหญ่แค่ไหนก็สู้ดูบนทีวีไม่ได้ ใต้ Wi-Fi เดียวกัน เปิดปุ่ม "TV" มุมขวาบนของแอปวิดีโอ ก็แคสต์ภาพขึ้นทีวีได้ ทีวีอัจฉริยะกับโปรเจกเตอร์ส่วนใหญ่รองรับ ดูซีรีส์ดูหนังสนุกขึ้นเยอะ ผู้ใช้ iPhone ยังใช้ AirPlay มิเรอร์ทั้งหน้าจอได้ด้วย
\end{itemize}

\booksection{สี่ อ่านหนังสืออิเล็กทรอนิกส์: บนมือถือก็อ่านได้อย่างสบาย}

มีคนว่าอ่านหนังสือบนมือถือทำร้ายตา และเสียสมาธิง่าย อันนี้จริงบางส่วน แต่จุดแข็งของการอ่านบนมือถือก็ชัด — อ่านได้ทุกที่ทุกเวลา ใช้เวลาเศษให้คุ้มที่สุด ฉันอ่านหนังสือไปไม่น้อยตอนอยู่บนรถไฟใต้ดิน ต่อแถว และก่อนนอน

\begin{itemize}
  \item \textbf{WeChat Reading}: ตอนนี้แอปอ่านที่ประสบการณ์ดีที่สุดในจีน ห้องสมุดครบ ทั้งหนังสือสำนักพิมพ์และนิยายออนไลน์ ใช้โมเดล "เวลาอ่านแลกเหรียญซื้อหนังสือ" อ่านมากก็แลกหนังสือฟรีได้ เป็นมิตรกับนักเรียนมาก ฟังก์ชัน "ขีดเส้นใต้" เห็นว่าคนอื่นขีดท่อนไหน บางทีเจอประโยคดีที่คนอื่นเมนต์ไว้ก็รู้สึกร่วมได้จริง แต่ "ตารางอันดับการอ่าน" อันนี้มันเว่อร์ไปหน่อย ง่ายกลายเป็น "อ่านเพื่อจาร์ยเวลา" แนะนำปิดทิ้ง
  \item \textbf{แอป Kindle}: ถ้าคุณมีเครื่อง Kindle แอป Kindle บนมือถือซิงก์ความคืบหน้าการอ่านให้ หนังสืออิเล็กทรอนิกส์ของแท้บน Kindle จัดหน้าดีกว่า PDF ปลอมเยอะมาก แต่ร้านหนังสือ Kindle โซนจีนปิดไปแล้ว ต้องสลับไปโซนต่างประเทศ
  \item \textbf{Apple Books}: มาในตัว iPhone เหมาะอ่าน EPUB กับ PDF อินเทอร์เฟซเรียบสุดๆ อนิเมชันพลิกหน้าสบายมือ ถ้าดาวน์โหลดไฟล์หนังสืออิเล็กทรอนิกส์จากเว็บมา "แชร์ไปที่ Books" ก็นำเข้าได้
  \item \textbf{คำแนะนำการตั้งค่าการอ่าน}: ไม่ว่าใช้แอปไหน ตั้งพื้นหลังเป็นโทนเหลืองอุ่น (โหมดปกป้องตา) ขยายฟอนต์ขึ้นหน่อย เว้นบรรทัดกว้างขึ้น อ่านตอนกลางคืนเปิด "โหมดมืด" — พื้นดำตัวขาว ตาสบายขึ้นเยอะ อย่าพึ่งอ่านพื้นขาวตัวดำตัวเล็กๆ ในที่แสงน้อยนานๆ ความเหนื่อยตามันสะสม
\end{itemize}

\booksection{ห้า เกมมือถือ: เล่นสนุก แต่ไม่ร้างเครื่อง}

ทัศนคติของฉันต่อเกมมือถือ: เล่นได้ แต่อย่าให้มันมาควบคุมคุณกลับ ใช้เครื่องมือเกมอย่างเหมาะสม เล่นสบายขึ้น งานสำคัญก็ไม่ตกหล่น

\begin{itemize}
  \item \textbf{โหมดเกม / พื้นที่เกม}: มือถือกระแสหลักตอนนี้มีกันหมด เปิดแล้วเครื่องจะเคลียร์พื้นหลัง บล็อกการแจ้งเตือน กันแตะพลาด — คุณกำลังรวมพรรคอยู่ สายด่วนโผล่มาครั้งเดียว โหมดเกมช่วยกันไว้ได้ บางเครื่องรองรับโหมดเพอร์ฟอร์แมนซ์/โหมดอีสปอร์ต บังคับชิปวิ่งเต็มสปีด แต่ราคาคือร้อนและกินไฟเร็วขึ้น
  \item \textbf{ปรับกราฟิกกับเฟรมเรตยังไง}: อย่าดึงเต็มสุดโดยไม่คิด เอารอบเฟรม (ความลื่น) ไว้ก่อน แล้วค่อยเรียกหากราฟิก 60 เฟรมสบายกว่า 30 เฟรมมากมาก แต่ 120 เฟรมกดดันมือถือมาก เครื่องระดับกลางไม่ต้องฝืน ถ้าเล่นเกมแข่งขัน (เช่นราชันยอดหยู หรือเกมยิงเอาชีวิตรอด) เฟรมเรตสำคัญกว่ากราฟิกมากๆ — จังหวะมือที่ลื่นสำคัญกว่าหญ้าสวยเป็นร้อยเท่า
  \item \textbf{มือถือร้อนตอนเล่นเกมทำยังไง}: ปกติ แต่ร้อนสูงนานๆ จะทำร้ายอายุแบต วิธีลดความร้อนมี: ถอดเคสให้ระบายความร้อน ลดความสว่างจอลงราว 70\% (ความสว่างคือแหล่งความร้อนตัวโต) อย่าชาร์จไปเล่นเกมไป (ความร้อนซ้อนกัน แบตพังไว) ถ้ามือถือร้อนถึงขั้นจับไม่ถนัดมือ หยุดให้มันพักก่อนเลย
  \item \textbf{เครื่องมือช่วยกันติดเกมมากไป}: "เวลาการใช้หน้าจอ" ของ iOS กับ "Digital Wellbeing" ของแอนดรอยด์ตั้งโควตาเวลาต่อวันให้แอปเกมได้ ฉันตั้งให้ตัวเองวันทำงานสูงสุดหนึ่งชั่วโมง สุดสัปดาห์สองชั่วโมง พอหมดเวลาไอคอนแอปจะเทา กดเข้าไม่ได้ ข้อจำกัดแบบนี้มีประโยชน์จริงๆ ไม่ใช่แค่เสี่ยงทาย
  \item \textbf{เกมคลาวด์}: ถ้าสเปกมือถือไม่พอแต่อยากเล่นเกมใหญ่ ลองแพลตฟอร์มคลาวด์เกม (เช่น Tencent START, NetEase Cloud Gaming) เกมรันบนเซิร์ฟเวอร์ระยะไกล ภาพส่งมาเป็นสตรีมวิดีโอ ต้องการเน็ตแรงแต่แทบไม่ขึ้นกับสเปกมือถือเลย ข้อเสียคือดีเลย์เลี่ยงไม่ได้ เหมาะเกมกลยุทธ์ RPG ไม่เหมาะเกมยิง FPS
\end{itemize}

\begin{tcbpractice}{ลงมือ: สร้างระบบความบันเทิงบนมือถือของคุณ}
\begin{enumerate}
  \item ดาวน์โหลดอัลบั้มหรือเพลย์ลิสต์โปรด 3 ชุดในแอปเพลง ให้แน่ใจว่าฟังออฟไลน์ได้
  \item ไปที่ Xiaoyuzhou หรือ Apple Podcasts กดติดตามช่องที่สนใจ 3 ช่อง ดาวน์โหลดหนึ่งตอนมาฟังดู
  \item ติดตั้ง VLC ลองใช้มันเล่นวิดีโอที่ก๊อปมาจากคอมสักไฟล์
  \item หาหนังสือที่อยากอ่านสักเล่มใน WeChat Reading ตั้งพื้นหลังสีอุ่นและโหมดปกป้องตา ขยายฟอนต์ให้เหมาะกับตัวเอง
  \item ถ้าคุณเล่นเกม เปิดโหมดเกมดูว่ามีการตั้งค่าไหนปรับได้บ้าง (กันรบกวน เฟรมเรตมาก่อน ฯลฯ)
\end{enumerate}

\end{tcbpractice}
